\documentclass[11pt,a4paper]{article}
\usepackage[utf8]{inputenc}
\usepackage[T1]{fontenc}
\usepackage{lmodern}
\usepackage[margin=24mm]{geometry}
\usepackage{amsmath,amssymb,mathtools}
\usepackage{booktabs}
\usepackage{enumitem}
\usepackage{xcolor}
\usepackage[hidelinks]{hyperref}

\newcommand{\dd}{\mathop{}\!\mathrm{d}}
\newcommand{\ii}{\mathrm{i}}
\newcommand{\id}{\mathrm{id}}
\newcommand{\R}{\mathbb{R}}
\newcommand{\Z}{\mathbb{Z}}
\newcommand{\N}{\mathbb{N}}
\newcommand{\Ord}{\operatorname{ord}}
\newcommand{\Ker}{\operatorname{Ker}}
\newcommand{\Aut}{\operatorname{Aut}}
\newcommand{\slashed}[1]{\not\!#1}
\setlength{\parindent}{0pt}
\setlength{\parskip}{6pt}
\allowdisplaybreaks

\title{Particle Physics --- Fermi's Golden Rule and the Dirac Equation}
\author{IMAPP lecture notes}
\date{Spring 2026}
\begin{document}\maketitle

\section{Fermi's golden rule}
For a perturbation $H_{\mathrm{int}}$ and an initial state $|i\rangle$, the first-order transition rate into final states with density $\rho(E_f)$ is
\[
\Gamma_{i\to f}=2\pi\,|T_{fi}|^2\rho(E_f),\qquad T_{fi}=\langle f|H_{\mathrm{int}}|i\rangle.
\]
For relativistic particle physics, one uses the invariant matrix element $\mathcal M$ and the invariant phase-space measure instead of the frame-dependent $T_{fi}$.

\section{The Dirac equation}
The Klein-Gordon equation describes scalar particles but does not provide a positive-definite probability density for a one-particle interpretation. Dirac sought an equation first order in time and space:
\[
\ii\partial_t\psi=\left(\boldsymbol\alpha\cdot\mathbf p+\beta m\right)\psi.
\]
Consistency with the relativistic dispersion relation requires
\[
\{\alpha_i,\alpha_j\}=2\delta_{ij},\qquad
\{\alpha_i,\beta\}=0,\qquad \beta^2=1.
\]
These matrices are represented by $4\times4$ Dirac matrices, so $\psi$ has four components.

Introducing $\gamma^0=\beta$, $\gamma^i=\beta\alpha_i$, the covariant equation is
\[
(\ii\gamma^\mu\partial_\mu-m)\psi=0.
\]

\section{Particle and antiparticle solutions}
The equation has positive- and negative-energy solutions. In the Feynman-Stueckelberg interpretation, negative-energy solutions are reinterpreted as positive-energy antiparticles propagating forward in time. The free spin-half solutions are written as spinors $u_s(p)$ for particles and $v_s(p)$ for antiparticles, with completeness relations of the form
\[
\sum_s u_s(p)\bar u_s(p)=\slashed p+m,qquad
\sum_s v_s(p)\bar v_s(p)=\slashed p-m.
\]

\section*{Summary}
The golden rule connects perturbative amplitudes to transition rates. The Dirac equation is Lorentz covariant, describes spin-$1/2$ fermions, and naturally predicts antiparticles.
\end{document}
